% Standalone author component. Requires amsmath, amssymb, amsthm, hyperref.
% All labels and new bibliography keys use q198bb:.
\providecommand{\BB}{\operatorname{BB}}
\providecommand{\Q}{\mathbb Q}
\providecommand{\Z}{\mathbb Z}
\section{A fixed 198-variable quartic selected by Busy Beaver}
\label{q198bb:specialization}

We retain the arithmetic hypotheses and the exact polynomial table of
 the 198-variable construction. The argument in this section bounds a
sufficient threshold by bounding the number of records in a computable
search. It then gives a separate strict bound on the coefficients.
Neither step evaluates a Busy Beaver value or computes an unknown
height constant by recognizing the last record.

\paragraph{Machine convention.}
Let $\BB(s)$ be the maximum blank-tape halting time of a deterministic
machine with exactly $s$ working states, a distinguished initial state,
symbols $0,1$, one two-way infinite tape initially all zero, and moves
left or right. A separate halting state is not counted. The transition
entering it, including its write and move, counts as one step. Unused
working states are allowed. The finite set of transition tables contains
a halting table, so the maximum exists. Padding proves that $\BB$ is
nondecreasing.

\begin{theorem}[The Busy Beaver specialization]\label{q198bb:fixed}
Subject to the arithmetic premises of the 198-variable construction,
let $Q_1(A),\ldots,Q_{184}(A)$ be its primitive integer quadratic
residuals in 198 total coordinates, with the first coordinate the input.
Use the unbalanced presentation, whose coefficients are polynomials in
$A$. Set
\[
 N_{\rm BB}=\BB(2^{1000}),\qquad
 A_{\rm BB}=149^{2N_{\rm BB}+1},\qquad
 P_{\rm BB}=\sum_{i=1}^{184}Q_i(A_{\rm BB})^2.
\]
Then, for every $x\in\Q$,
\[
 x\in\Z\quad\Longleftrightarrow\quad
 \exists(z_2,\ldots,z_{198})\in\Q^{197}:
 P_{\rm BB}(x,z_2,\ldots,z_{198})=0.
\]
The polynomial has degree four, content one and exactly 2945 nonzero
coordinate monomials. Its largest absolute coefficient is
\[
 H(P_{\rm BB})=149^{4\BB(2^{1000})+2}<\BB(2^{2000}).
\]
The parameter $A_{\rm BB}$ is fixed and is not an additional quantified
coordinate. The count 198 is an upper bound, not a minimum-variable claim.
\end{theorem}

\subsection{The arithmetic inputs and their precise scope}
\label{q198bb:inputs}
We state the conclusions supplied by the preceding arithmetic proof
that are used below. Write $h$ for absolute logarithmic Weil height and
$\widehat h_E$ for the elliptic canonical height.
Put $K=\Q(\rho)$, $\rho^2=2$, $d=3-\rho$, $D=d^2$ and
\[
 E:y^2=x^3-Dx,\qquad R_0=(1+\rho,2\rho),\qquad
 P=2R_0=\left((57-28\rho)/8,(-660+417\rho)/32\right).
\]
The retained group and height statements give
$E(K)=\Z R_0+E[2](K)$, uniquely, with
$E[2](K)=\{O,(0,0),(d,0),(-d,0)\}$ and
$\widehat h_E(P)\ge1/2$. The two pointed images of $R_0$ on
$E_0:Y^2=X^3-X$, obtained from the conjugate curves by their fixed
pointed isomorphisms, are independent modulo torsion over $\Z[i]$.
The fixed real-orbit recurrence, finite-place preparation, compatible
CRT, simultaneous-prime and prepared pole-realization statements are
those of the 198 construction.

Every rational solution recovers four finite points
$Q_i=[m_i]R_0+\tau_i$, with $m_i\in\Z$ and $\tau_i\in E[2](K)$,
and labels
\[
 n=m_1/2,\qquad a=m_2\ne0,\qquad t=m_3,\qquad b=2m_4+1.
\]
Its aggregate and scalar data are
\[
 T=7+\frac14\operatorname{Tr}_{K/\Q}\sum_i x(Q_i)\ge1,
 \quad t_0=\lambda,\quad t_1=\lambda+T,
 \quad B=\max(1,|n|,|a|,|t|,|b|),\quad c=|a|,
\]
\[
 \beta=51x^2+2504,\qquad
 L_N=\lfloor9(2N-1)/4\rfloor+1\ge N,
\]
\[
 c\ge149^{L_N},\qquad
 h(\beta)\le C_3 149^{-L_N}c(1+\log(c)/3),\qquad C_3=2^{330770}.
\]
The retained contact radicals $M_j=M(t_j)$ satisfy
\[
 h(t_j)\le H_2M_j^{7/6},\qquad
 M_j^3\mid\Delta\quad(\Delta\ne0),\qquad
 M_j\mid J_x\quad(J_x\ne0),
\]
where
\[
 \Delta=t^2-(a+1)(a-1)tb+(a-1)^2b^2-1,\quad |\Delta|\le14B^4,
 \qquad J_x=|\operatorname{num}(\beta-n)|\le4B e^{h(\beta)}.
\]
Here $\operatorname{num}$ means reduced numerator, and the divisibility
statements concern nonzero integers. The contact supplier has
$C_2<2^{7557/32}$ and $\alpha=2103/1024<7/3$, with
$H_2=\log7+C_2B_\Sigma^\alpha$ and
$B_\Sigma=\prod_{p\in\Sigma\setminus S_{10}}p$. The retained exceptional
set consists of primes at most 76176 together with 2029147.
Consequently
\[
 B_\Sigma\le2029147\,76176^{76176}<2^{1295013},\qquad
 H_2<2^{4000000}.
\]
The norm-one/Lucas classification gives the following explicit
consequence when $\Delta=0$ and $c\ge8$. Put $A'=a+1$,
$c'=|A'|$ and $w=\operatorname{sgn}(A')(A'-2)b$. Then $w\ne0$ and
$t^2-c'tw+w^2=1$. The retained classification gives
$|w|=U_j(c')$, $|t|\in\{U_{j-1}(c'),U_{j+1}(c')\}$ and
$j\ge c'-2$. The recurrence bound
$U_\ell(c')\ge(c'-1)^{\ell-1}$ for $\ell\ge1$ therefore gives
\[
 \log B\ge(c'-5)\log(c'-2)\ge(c-6)\log(c-3).
\]
All bases and exponents in this weakening are positive for $c\ge8$.
When $J_x=0$, the reduced-denominator argument using squarefreeness of
102 gives $x\in\Z$.

For completeness the actual scalar quotient is
$Z=2z(Q_2)/z(H)$, where $z=-x/y$, and
\[
 R=\frac{149^{2N+1}r_s}{1+149^2r_s},\quad
 r_s=\beta(\beta-1)(\beta-2401)(32\beta+343)\ge1,
 \quad F=\operatorname{Nm}_{K/\Q}(1/x(H)),
\]
\[
 W_s=Z_0^4/R^{10}+Z_1^2/R^5+F^2/R^4,\qquad
 b_s=R(R+9792)W_s^2.
\]
The ordered completeness construction permits arbitrary required
finite depths for $H$ and then $Q_2=[a_0]P$, including $3\mid a_0$,
while keeping their fixed real neighborhoods. After their values and the
labels have been chosen, the comparison-prime theorem applies to the
fixed rational contact list, merging repetitions and retaining incidences.
The final progression preserves the earlier preparation. Prepared pole
realization is applied to each actual final input after $H$ is fixed.
The only uniform lower restriction on $N$ in this construction is the
one ensuring $0<b_s<17$; the later moduli, primes, slopes and certificates
may depend on $(N,x)$ and impose no new lower bound on $N$.
Section~\ref{q198bb:real} supplies an explicit sufficient uniform bound.

Finally, the exact unbalanced coefficient calculation gives, for $N\ge1$,
\[
 H(P_N)=\max\{H_0,149^{4N+2}\},\quad
 K_0=31581660^{28},\quad
 H_0=(5898050446009651968K_0)^2<149^{214}.
\]
The degree, content and 2945-monomial support hold throughout this
parameter range. The separate balanced table uses $C=149^N$ and has
coefficients linear in $C$; it is not the polynomial-in-$A$ presentation
selected in Theorem~\ref{q198bb:fixed}.

\subsection{The enlarged point domain and the fixed geometry}
\label{q198bb:geometry}
For height purposes replace each $m_i$ by $|m_i|$: this preserves its
abscissa since $-\tau_i=\tau_i$. Write
$S=\sum_i m_i^2$. Freeze every zero-index coordinate to one of the three
finite two-torsion points, and omit that factor. There are at most
$4^4=256$ strata: in each position choose an active index or one of three
fixed torsion points. Active torsion translates remain unrestricted.
An all-frozen stratum has $S=0$ and produces no record below.

For $k\ge1$ active positions use $A_k=E_0^{2k}$, embedded by the product
of the complete plane-cubic embeddings and Segre. The cubic exact
sequence gives surjectivity of restriction of homogeneous forms in
every nonnegative degree; tensoring and Segre give the corresponding
condition (1.1) of \cite{q198bb:VB}. The line bundle is ample and
symmetric. On $A_k$ use the function
\[
 F=c_0+\frac d4\sum_{i=1}^kX_i+\frac {d'}4\sum_{i=1}^kX_i'.
\]
The constant $c_0$ includes 7 and the omitted contributions. Each omitted
contribution is $0$ or $\pm3/2$, so $1\le c_0\le13$, its denominator
divides 2, and $h(c_0)\le\log26$.

We recall why the actual tuples avoid every exceptional abelian
translate. A positive-dimensional abelian subvariety of a power of
$E_0$ contains an elliptic subvariety isogenous to $E_0$. A direction
through the tuple therefore has coordinate endomorphisms
$\gamma_i,\gamma_i'\in\Z[i]$, not all zero. If two nonconstant
first-block coordinates share a pole, elimination of its parameter gives
$m_i\gamma_j=m_j\gamma_i$. The coefficient of their common double pole,
including all terms with that pole, is proportional to
\[
 \frac d{4\gamma_i^2}\sum_j(m_i/m_j)^2\ne0.
\]
It cannot cancel. A pole shared across the two conjugate blocks would
give a nonzero relation between the pointed images of $R_0$ modulo
torsion, contrary to their retained independence. The same positive
sum works in the second block. Thus $F$ is nonconstant on any such
direction. The graph closure $V\subset A_k\times\mathbf P^1$ is
integral; a fiberwise translate through an actual domain point meets
the domain densely, so graph boundary components do not evade this
argument.

Multiplying its graph equation by the product of the $2k$ homogeneous
$Z$ coordinates gives a hyperplane section in the Segre embedding of
$A_k\times\mathbf P^1$. The graph is an irreducible component, whence
\begin{gather*}
 \dim V=2k\le8,\qquad m=2k+1\le9,
 \qquad N_A+1=3^{2k}\le3^8,\qquad N_P+1=2,\\
 D_V\le(2k+1)!3^{2k}\le9!3^8=2380855680<2^{32}.
\end{gather*}
The paired canonical height is
\[
 \widehat h_{A_k}=\alpha S,\qquad
 \alpha=6\widehat h_E(R_0)=\tfrac32\widehat h_E(P)\ge\tfrac34.
\]
Torsion does not affect it. Normalized directions have exactly the
Euclidean metric of $(m_i)/\sqrt S$, with zeros in omitted positions,
in a four-dimensional cube. This covers the full point domain needed here.


\subsection{Small fixed height constants.}
We record generous bounds in the height convention of the projective
embedding. The duplication map on the abscissa of $E_0$ is
\[
 (X:Z)\longmapsto\big((X^2+Z^2)^2:4XZ(X^2-Z^2)\big).
\]
Writing $F=(X^2+Z^2)^2$ and $G=4XZ(X^2-Z^2)$, the identities are
\[
\begin{aligned}
 Z^7&=(Z^3-\tfrac34X^2Z)F
              +(\tfrac5{16}XZ^2+\tfrac3{16}X^3)G,\\
 X^7&=(X^3-\tfrac34XZ^2)F
              -(\tfrac3{16}Z^3+\tfrac5{16}X^2Z)G.
\end{aligned}
\]
Each row has sum of absolute coefficients $9/4$. The local height estimates give
\[
 4h(x)-\log36\le h(x([2]Q))\le4h(x)+\log8.
\]
The estimate includes projective abscissas of torsion points. Telescoping
gives $|h(x(Q))-2\widehat h_E(Q)|\le\log36/3$.
The plane height differs from $\tfrac32h(x(Q))$ by at most $\log4$,
as follows directly from $y^2=x^3-x$ at each place. It follows that
$C_{\rm haut}<100$ on $A_k$ with the usual Weil height on the base.

For addition and subtraction, the nine products of the coordinates of
$Q+R$ and $Q-R$, multiplied in the affine chart by $(x_R-x_Q)^3$,
reduce to polynomials of bidegree $(2,2)$ with integer coefficients of
absolute value at most four. The nine polynomials are displayed in
Section~\ref{q198bb:finite-algebra}. Their extensions are sections of
$\mathcal O(6O)\boxtimes\mathcal O(6O)$, by the theorem of the square.
The divisor of $x_R-x_Q$ identifies the common trivialization used in
this calculation, so these sections give the addition-subtraction
morphism also across the exceptional affine charts.
Taking at most eight products and expressing them in Segre coordinates gives
$h_{\rm add}<100$; even bounding every coefficient by $2^{40}$ suffices.


For the successive-minima step we use the final published
\cite[Theorem~5.2, p.~208]{q198bb:Zhang}, rather than relying on
a comparison between the author text of \emph{Small points and adelic
metrics} and its published edition.  Take the projective integral
closure of the graph over the ring of integers of a field of definition,
and the restriction of $\mathcal O(1)$ with its standard Fubini--Study
metric.  The generic line bundle is ample.  Its finite-place model is
relatively generated and its smooth metric is the projective quotient
metric, so it is relatively semiample in the terminology of that theorem.
The arithmetic dimension is $2k+1\le9$.  The theorem's upper inequality is
\[
 \frac{h_{\rm FS}(V)}{(2k+1)\deg V}\le e_1(V).
\]
Here the arithmetic intersection height is divided by the field degree;
the same field degree occurs in the geometric degree, so no extension
factor remains.

Torsion tuples outside the poles of $F$ are Zariski dense in $A_k$.
Their graph points are therefore dense in $V$.
The duplication bound gives abscissa height at most $\log36/3$ on
each factor and plane height less than $10$.  The at-most-eight-term height inequality, together with
$h(d/4),h(d'/4)\le\tfrac12\log7+\log4$, bounds the active trace part
by $40$. Adding $c_0$ costs at most $\log26+\log2<4$, so their
$F$ values have height below $44$. Their graph points have Weil height
below $124$. The Euclidean metric adds at most
$\tfrac12\log(2\cdot3^8)<5$. In particular $e_1(V)\le200$, where
this $e_1(V)$ denotes the essential minimum, not the later Vojta constant.
The published inequality yields
$h_{\rm FS}(V)\le(2k+1)D_V\cdot200\le9D_V\cdot200$.

For clarity, the coarse Chow-height comparison used here can be made
with $m=2k+1\le9$ blocks, $n=2\cdot3^{2k}\le2\cdot3^8$ coefficients per block, and block
degree $D_V$.  Repeated one-variable coefficient--Mahler inequalities
give $\log\|\operatorname{coeff}Q\|_\infty
\le m_{\rm tor}(Q)+mnD_V\log2$.
For sphere measure, the mean logarithm of a coordinate radius is
$-H_{n-1}/2$.  Convexity of the torus logarithmic mean in the logarithms
of the radii and homogeneity give
$m_{\rm tor}(Q)\le m_{\rm sph}(Q)+mD_VH_{n-1}/2$.
Here is the needed intersection calculation, so that no uninspected
published comparison formula is required.  Write $r=\dim V=m-1$,
$D=\deg V$, and $H_{n-1}=\sum_{k=1}^{n-1}1/k$.  For a unit vector
$u\in\mathbb C^n$ put
\[
 \ell_u([z])=\langle u,z\rangle,\qquad
 \log\|\ell_u([z])\|=\log|\langle u,z\rangle|-\log\|z\|_2.
\]
Let $u_0,\ldots,u_r$ be independent unit vectors with invariant
probability measure.  Outside a set of measure zero their successive
hyperplanes meet $V$ properly.  Set
$V_j=V\cap\{\ell_{u_0}=\cdots=\ell_{u_{j-1}}=0\}$, with
$V_0=V$, and let $\omega=c_1(\mathcal O(1),\|\cdot\|_{\rm FS})$.
Every $V_j$ has degree $D$.  The recursive definition of a metrized
intersection (or, equivalently, integration by parts with
$dd^c\log\|\ell_u\|=[\ell_u=0]-\omega$) says that replacing
one hyperplane section by its zero divisor changes the logarithm
of its intersection symbol by
\[
 I_j=\int_{V_j(\mathbb C)}\log\|\ell_{u_j}\|\,\omega^{r-j}.
\]
Applying this recursion for $j=0,\ldots,r$ gives the local Chow
symbol: its zero divisor in the hyperplane variables is precisely
the locus where the $r+1$ hyperplanes have a common point on $V$,
with the intersection multiplicities.  Thus its polynomial is the
Chow form $Q$, up to a scalar.  In this identification the recursion
reads
\[
 \log|Q(u_0,\ldots,u_r)|
       =-\log\|q\|_{\rm int}+\sum_{j=0}^r I_j,
\]
where $q$ is the corresponding basis of the Chow line and
$\|q\|_{\rm int}$ is its intersection metric.  The sign can also be
checked in dimension zero: for $Q(u)=\langle u,z\rangle$ the right
side is $\log\|z\|_2+\log\|\ell_u([z])\|$.
This argument uses intersections of cycles and integration over their
regular loci, so it does not require $V$ to be smooth.

For fixed $[z]$, unitary invariance identifies
$|\langle u,z\rangle|^2/\|z\|_2^2$ with the first squared coordinate
of a uniform unit vector.  Its density on $[0,1]$ is
$(n-1)(1-t)^{n-2}$, whence
\[
 \mathbb E_u\log\|\ell_u([z])\|
 =\tfrac12(n-1)\int_0^1\log t\,(1-t)^{n-2}\,dt
 =-\tfrac12 H_{n-1}.
\]
Conditioning on $u_0,\ldots,u_{j-1}$ therefore gives
$\mathbb E I_j=-DH_{n-1}/2$.  These logarithms are integrable;
indeed their negatives have the displayed finite integral, so
Tonelli's theorem justifies the successive averages.  The archimedean
intersection contribution is consequently
$m_{\rm sph}(Q)+mDH_{n-1}/2$.

At a finite place the same recursion uses the standard model
$\mathcal O(1)$.  Take the hyperplane coefficients to be independent
indeterminates with Gauss norm one.  A polynomial evaluated there
has norm equal to its coefficient Gauss norm; the elimination symbol
just identified with $Q$ therefore contributes
$\log\|Q\|_{v,\rm Gauss}$.  This also follows by taking a primitive
local Chow form: its reduction is nonzero, and generic hyperplanes
over the residue function field avoid its zero locus.  Arbitrary
scaling of $Q$ changes these local terms by $\log|c|_v$, whose weighted
sum is zero.  Summing the local recursion and dividing by the field
degree proves
\[
 h_{\rm FS}(V)=\frac1{[K:\mathbb Q]}
 \left(\sum_{v\nmid\infty}[K_v:\mathbb Q_v]
       \log\|Q\|_{v,\rm Gauss}
 +\sum_{\sigma:K\hookrightarrow\mathbb C}
       m_{\rm sph}(Q^\sigma)\right)+\frac{mD}{2}H_{n-1}.
\]
This supplies directly the comparison needed here.  The published
Chow-form discussion in \cite[Section~4.3.4.1]{q198bb:BGS} is background;
the displayed identity has been derived by the local intersection recursion.
Thus the larger error
$mnD_V\log n$ bounds the differences needed for the coefficient and
unitary Chow conventions.  Here this is
$9D_V(2\cdot3^8)\log(2\cdot3^8)<2^{60}$.
Thus $h(V)<2^{100}$ in the convention consumed by the family theorem.
Together with the displayed finite addition and duplication calculations,
this preserves the estimates $h_{\rm add},C_{\rm haut}<100$ and all
the record-count constants $\Lambda,\psi,e_i,\mu,U,K$ used below.




\subsection{A finite number of computable records}
\label{q198bb:records}
Apply the quantitative theorem \cite[Theorem~3.1, p.~312]{q198bb:VB}
to each stratum. Its constants are given by (3.2) in that source:
\[
 \Lambda=2^{m^2}(8m\dim V)^{m\dim V}(N_A+1)(N_P+1)D_V^{3m},
 \qquad \psi=\prod_{j=2}^{m\dim V}(3j+1),
\]
\[
 e_1=8\Lambda^\psi,\quad e_2=2\Lambda^{2\psi},\quad
 \mu=4\Lambda^\psi,\quad
 e_3=(4(N_A+1))^{2m^2\dim V}\Lambda^{8\psi}
                  \max(h(V),h_{\rm add},C_{\rm haut}).
\]
It forbids $m$ points outside the fiberwise translate locus when
consecutive normalized directions have distance at most $1/e_1$,
consecutive heights increase by at least $e_2^2$, the first height is
at least $e_3$, and every height is at least $\mu$ times its base height.
We may majorize all constants by evaluating the displayed formulas at
the bounds for $k=4$. Thus we use common upper envelopes below, and
nine points suffice to contradict the theorem in any stratum.

The geometry estimates give
\[
 \Lambda<2^{1700},\qquad \psi<2^{568},\qquad U=2^{2^{600}}.
\]
In particular
\[
 \log_2 e_3<19440+13600\cdot2^{568}+100<2^{583},
 \qquad \mu,e_2^2,e_3<U/4.
\]
Cover $[-1,1]^4$ by cubes of side at most $1/(2e_1)$, with at most
$\lceil4e_1+1\rceil^4$ cells and diameter at most $1/e_1$.
Each cell in each stratum can contain at most eight records satisfying
the height conditions. Put $K=8\lceil4e_1+1\rceil^4$.
The bound $\log_2K<2^{584}$ gives $256K<U$.

For a rational $T=u/v$ in lowest terms with $v>0$, put
\[
 b(T)=\operatorname{bitlength}(\max(|u|,v));\qquad
 h(T)\le b(T)\le2h(T)+1.
\]
Encode a finite elliptic point by a positive integer $q=4m+e$, where
$m\ge0$ and $e=0,1,2,3$ denotes the torsion translate
$O,(0,0),(d,0),(-d,0)$ respectively. The only excluded pair is $(m,e)=(0,0)$.
Enumerate all positive code quadruples by their maximum code and then
lexicographically within each finite shell. Compute their exact rational
aggregate and $S=\sum_i\lfloor q_i/4\rfloor^2$.

Initialize $r=U$. Declare a record exactly when
\[
 S>Ur\quad\hbox{and}\quad S>Ub(T),
\]
and then replace $r$ by $S$. These predicates use only exact integer and
quadratic-field arithmetic. They require no comparison of computable
real numbers. Each recorded canonical height is at least $3S/4$, hence
exceeds $e_3$ and $\mu h(T)$. Consecutive records have height ratio
exceeding $U>e_2^2$. A subsequence in a fixed stratum and direction cell
preserves this growth. It follows that the number $j_*$ of records is
finite and $j_*<256K<U$.

Let $r_*=U$ if there are no records, and otherwise let $r_*$ be the last
record. Every tuple satisfying $S>Ub(T)$ has $S\le Ur_*$: a later
violation would be another record, and every earlier running value is
at most $r_*$. Consequently
\[
 S\le Ub(T)+Ur_*\le2Uh(T)+U(r_*+1).
\]
For the signed labels of Section~\ref{q198bb:inputs},
$|2m_4+1|\le2|m_4|+1$ gives
\[
 B^2\le9(S+1)\le18Uh(T)+10U(r_*+1).
\]
The last inequality uses $U(r_*+1)\ge9$. Thus the uniform aggregate
constants $M_1=18U$ and $M_2=10U(r_*+1)$ cover every rational solution,
including zero-index torsion differences.

\subsection{The all-larger-parameter arithmetic threshold}
\label{q198bb:threshold-arithmetic}
Set
\[
 G(r)=U^{32}(r+1),\qquad N\ge G(r_*),\qquad
 C=2M_1H_2=36UH_2<U^2,\qquad J_1=32U(r_*+1)<G(r_*).
\]
The extra $M_1\log2$ in
$h(T)\le h(t_0)+h(t_1)+\log2$ is absorbed by $J_1$.
The retained scalar inequalities give $B\ge c\ge149^{L_N}>N$.
If $\Delta\ne0$, cubic contact divisibility gives
\[
 B^2\le C14^{7/18}B^{14/9}+J_1
      \le C14^{7/16}B^{7/4}+J_1.
\]
Since $C14^{7/16}<4U^2$ and $B>U^{32}$, the first term is less than
$B^2/2$. Since $J_1<N<B$ and $B>2$, so is the second. This contradiction
proves $\Delta=0$.

If $J_x\ne0$, its size bound and linear contact divisibility, after
absorbing $J_1<B^2/2$, give
\[
 \log B\le\tfrac65\log(2C)+\tfrac75\log4+\tfrac75h(\beta)
 \le C'+7C_3 149^{-L_N}c(1+\log(c)/3),
\]
where $C'=4\log(2C)+7\log2<U$. Our lower bound on $N$ implies
$7C_3 149^{-L_N}<1$ and $c>\max(C',e^{12})$.
The Lucas inequality gives
\[
 \log B\ge(c-6)\log(c-3)>2c+c\log(c)/3
                         >C'+c+c\log(c)/3,
\]
contradicting the upper bound. For the strict middle inequality one may
use $c-6>0.9c$ and $\log(c-3)>\log c-1$ when $c\ge e^{12}$.
Thus $J_x=0$ and $x\in\Z$. All inequalities hold for every larger $N$;
no monotonicity assertion about the zero sets themselves is used.

\subsection{Uniform explicit real preparation.}\label{q198bb:real}
We give an explicit bound independent of the input and of all later
finite moduli.  This avoids assuming that a particular multiple such
as $64P$ lies in the identity component of the closure of the paired
real orbit.

For $\epsilon\in\{+1,-1\}$ put $d_\epsilon=3+\epsilon\sqrt2$ and
$E_\epsilon:y^2=x^3-d_\epsilon^2x$.
Let $P_\epsilon$ be the image of $P$ under
$\sqrt2\mapsto-\epsilon\sqrt2$.
Let $G$ be the closure of the cyclic subgroup generated by
$(P_+,P_-)$ in $E_+(\mathbb R)\times E_-(\mathbb R)$, and let
$G^0$ be its identity component.
Both projections of $G^0$ onto the real identity circles are surjective.
Indeed each projected $P_\epsilon$ has infinite order.  The image of
$G$ under either projection is therefore infinite.  If the image of
$G^0$ were trivial, that image would be a quotient of the finite
component group $G/G^0$, a contradiction.  The image of $G^0$ is a
connected closed subgroup of a circle, so it must be the whole circle.
Probability Haar measure on $G^0$ therefore pushes forward to
probability Haar measure on each circle.  No independence assertion
about the two projections is needed.

For $1<d<5$, the identity component of $E_d$ has finite abscissa
$x\ge d$.  Its normalized invariant measure is obtained from
$dx/(2|y|)$ on the two branches.  Write
\[
 \Omega_d=\int_d^\infty \frac{dx}{\sqrt{x^3-d^2x}}.
\]
Integration on $[d,2d]$ gives
$\Omega_d\ge1/(2\sqrt2\sqrt d)$.
For $X\ge2d$, integration on $[X,\infty)$ gives
\[
 \int_X^\infty \frac{dx}{\sqrt{x^3-d^2x}}
           \le \frac{4}{\sqrt3\sqrt X}.
\]
Moreover $|z|^2=x/(x^2-d^2)>1/x$ for $z=-x/y$.
Thus the Haar measure of $|z|<1/1000$ is less than
$8\sqrt{2d}/(1000\sqrt3)<32/1000$.
The union bound shows that a point of $G^0$ exists where both real
$z$ values have absolute value greater than $1/1000$.
The finite exceptional points and the equality level sets have measure
zero, so they can be avoided.  Fix a compact neighborhood of such a
point, within $G^0$, on which both $|z|>1/2000$ and the reference
charts are defined.

Let $m=|G/G^0|$.  Enlarge any finite preparation modulus to a multiple
of $m$ and $64$.  Its multiples of the paired point are dense in
$G^0$: multiplication by that modulus kills the component group and
is surjective on the torus $G^0$.
The retained positive-orbit recurrence therefore chooses the initial
reference $H$ in the fixed neighborhood with all required finite
divisibilities.  The scalar point $[a_0]P$ is then chosen in a fixed
neighborhood of the identity with both $|z([a_0]P)|\le1$, and with the additional fixed divisibility $3\mid a_0$.
Continuity fixes a residue class for the later reference choices,
and the retained CRT and real-progression argument preserves it.
This is the same order of choices as the 198-variable completeness
argument; no bound on $m$ or on the final indices is required.

The two real values of the actual quotient
$Z=2z([a_0]P)/z(H)$ have absolute value below $4000$,
so its rational components satisfy $|Z_0|,|Z_1|<2^{12}$.
Also $|F|=|\operatorname{Nm}(1/x(H))|\le1/7$, since the two positive
real abscissae are at least $d_+,d_-$ and $d_+d_-=7$.
The actual scalar expression (writing $R_s=R$) is
\[
 W_s=\frac{Z_0^4}{R_s^{10}}+\frac{Z_1^2}{R_s^5}+\frac{F^2}{R_s^4},
 \qquad b_s=R_s(R_s+9792)W_s^2.
\]
For $R_s\ge1$ it satisfies
\[
 0<W_s<2^{49}R_s^{-4},\qquad 0<b_s<2^{112}R_s^{-6}.
\]
In particular $R_s\ge2^{20}$ suffices for $0<b_s<17$.
The uniform lower bound
$R_s\ge149^{2N+1}/(1+149^2)$ and $N\ge G(r_*)$ are more than enough.
The factor two is a unit at every odd preparation prime and increases
the required $2$-adic depths. The finite preparation can still be made
arbitrarily deep. All subsequent finite, CRT, contact-prime, slope and
pole-certificate choices may depend on $(N,x)$ and impose no further
lower bound on $N$, by the actual 198 completeness construction stated
in Section~\ref{q198bb:inputs}. All labels and $T$ are fixed before
comparison primes; repeated contacts are merged with their incidences.
The pole predicate is applied to the final scalar input after $H$ has
been fixed. Thus every $N\ge G(r_*)$ gives completeness for every
integer input. Together with Section~\ref{q198bb:threshold-arithmetic}
this proves that $G(r_*)$ is an all-larger-parameter defining threshold.
In particular $G(r_*)>53$.




\subsection{The halting selector and its threshold}
\begin{theorem}\label{q198bb:threshold}
There is a sufficient defining threshold $N_*<\BB(2^{1000})$ for the
198-variable family. Every $N\ge N_*$ defines the integers.
\end{theorem}
For each integer $0\le j<U$, the program of
Section~\ref{q198bb:programs} enumerates the encoded finite point
quadruples until its $j$th record. When $j=0$ it immediately uses $r=U$.
It then computes $G(r)=U^{32}(r+1)$, performs $G(r)$ increments, and halts.
For $j=j_*$ it halts: the final record occurs in a finite shell, all
preceding shells are finite, and every exact arithmetic call terminates.
Programs with an input larger than the actual record count may run
forever. No instruction attempts to decide that a record is final.
The following compilation charges the hardwired finite ordinal $j_*$;
it does not hardwire the unknown value $r_*$ or the threshold.

\subsection{A constructive finite-state budget}\label{q198bb:compiler}
We use deterministic one-tape machines on a two-way infinite binary
tape, initially all zero, with moves left or right. The initial
working state is distinguished, the halting state is not counted,
and the transition entering it counts as a step. Unused working
states are permitted. These conventions fix the function $\BB$ used
below.

Here is a compilation bound for the two displayed integer programs.
It concerns their mathematical unbounded-integer semantics, not the
resources available to a particular Python interpreter. We first
compile to counter instructions
\[
 \operatorname{INC}(i,a),\qquad
 \operatorname{DECJZ}(i,a,b),\qquad \operatorname{HALT}.
\]
The second instruction decrements counter $i$ and goes to $a$ when
that counter is nonzero; otherwise it goes to $b$. All counters are
natural numbers. A dedicated counter kept zero implements an
unconditional jump by a single zero test.

For clarity, the following bounds include all loop and exit labels.
Clearing a counter uses one decrement loop. Moving a counter into
another uses a decrement followed by an increment. To copy $a$ into
$b$ while preserving $a$, first clear $b$ and a scratch counter $t$,
move one unit at a time from $a$ to both $b$ and $t$, then move $t$
back into $a$. This uses at most ten instructions. To compare two
natural numbers, copy them and decrement the copies in parallel,
recording which becomes zero first. Truncated subtraction uses the
same loop and retains the remaining first copy. Multiplication
copies the second factor to a loop counter and repeatedly adds the
first factor to an initially zero result. Division with remainder
by a positive number repeatedly compares and subtracts that number
from a copy of the dividend, incrementing the quotient each time.
These constructions give the following static instruction bounds;
their running times need not be bounded here:
\[
\begin{array}{c|rrrrrrr}
\text{operation}&\text{clear}&\text{move}&\text{copy}&\text{add}
 &\text{compare/subtract}&\text{multiply}&\text{divide/remainder}\\
\hline
\text{instructions}&1&3&10&24&64&64&256.
\end{array}
\]
For example, division contains one comparison template, one
subtraction template, initial copies and the quotient increment;
it contains no unrolled subtraction loop. Every scratch counter is
cleared before use, and no bound depends on the values in counters.
The input counters of copy, addition, comparison, multiplication and
division are preserved; their output counters are distinct from
their inputs and from each other. Move is explicitly destructive.
Its source and destination are distinct. Inputs may coincide with
each other. All generated scratch names belong to a reserved
namespace disjoint from source variables and live arguments.
Truncated subtraction has both a fresh-output version and a version
that updates its first argument; division uses the latter only
after a successful comparison. Before overwriting a source-language
variable, expression evaluation writes into fresh temporary
counters, then copies the result into that variable. Scratch
counters belonging to distinct expanded nodes are disjoint. These
contracts also apply to signed magnitudes and signs, tuple fields,
\texttt{None} tags and return-value copies. Fields not used when a
tag is \texttt{None} are initialized to zero and are not inspected
as tuple values on that branch.

Represent a signed integer by a sign bit and a magnitude, with sign
zero for magnitude zero. Addition and comparison split into the
four sign cases and use the preceding natural-number operations.
For division, first divide the magnitudes. In the opposite-sign
case with nonzero remainder, increase the quotient magnitude by
one and replace the remainder magnitude by the divisor magnitude
minus the old remainder. Assign the remainder the divisor's sign.
This implements floor division and the corresponding remainder,
including negative arguments. Fewer than $2^{12}$ counter
instructions, and fewer than that many scratch counters, suffice
for each signed operation: even four separate copies of the
256-instruction division template, the adjustments, copies and
sign tests fit this bound. A literal of magnitude at most 1000
uses at most 1001 instructions. Conditions use these comparisons
and jumps; while loops branch back to their condition. Boolean
short-circuiting is preserved by branches.

The record program has 1398 syntax nodes, fifteen functions and
acyclic call depth ten. Each function has fewer than $2^{12}$ nodes.
Using the looser depth bound twelve, the inlining recurrence
$a_d\le2^{12}(1+a_{d-1})$ is bounded by $2^{13d}$, giving fewer than
$2^{156}$ fully inlined nodes. The coefficient interpreter has 302
nodes and call depth three; the bounds $2^{11}$ per function and depth
nine give fewer than $2^{110}$ inlined nodes for it. Fixed tuples have
length at most three and nesting depth at most two. Flattening them needs at most nine
integer fields and a tag for \texttt{None}; returns become copies
to designated fields followed by a jump. Componentwise equality,
tuple unpacking, tags and return handling expand each node into
fewer than $2^{10}$ scalar operations. Thus fewer than $2^{166}$
scalar nodes suffice for the record program and fewer than $2^{128}$
for the coefficient interpreter. The only executed primitives are
the operations just described, tuple-field selection at fixed
indices, assignments, tests and loops. There are no recursive,
library or dynamically dispatched calls.

Allowing $2^{20}$ counter instructions and fresh counters per
expanded node is more than the explicit bounds above. Thus the
record core has fewer than $2^{186}$ instructions and counters;
the coefficient core has fewer than $2^{148}$. All divisions in the
record search have nonzero divisors: the field routines normalize
genuine elements of $\mathbb Q(\sqrt2)$, every positive code describes
a finite point, and the chord/tangent routine handles the identity,
inverse and two-torsion cases before dividing. The coefficient interpreter divides only by positive
powers of two. Consequently the arithmetic templates terminate
on every invocation that occurs in a halting computation.

We next give the counter-to-tape bound. A machine with $L$
counter instructions and $R$ counters represents its data as
\[
 \mathtt{L}\,\mathtt{\#}\,1^{r_0}\,
 \mathtt{\#}\,1^{r_1}\cdots
 \mathtt{\#}\,1^{r_{R-1}}\,\mathtt{E}
\]
on a tape with blank symbol distinct from these four symbols.
At every instruction boundary the head is at $\mathtt{L}$.
To access counter $i$, the finite control counts at most $R$
separators while scanning right. Increment inserts a mark before
the counter's next separator: carry that separator right, then
carry each displaced symbol right, until the end marker has been
moved onto the next blank. Decrement deletes the first mark by
copying each following symbol one cell left through the end marker
and erasing its former cell. An empty block takes the zero branch.
Each shifting routine stores just one of the five tape symbols
and a phase in finite control. It uses fewer than 32 states.
Returning left to $\mathtt{L}$ takes a separate scan for each
continuation. Including the initial writing of the $R$ empty
blocks, all locating, shifting, branching and return states fit
within $128(L+1)(R+1)$ states over at most eight symbols.

Encode each tape symbol by a three-bit word, choosing 000 for
blank. At a simulated step boundary the head is at the first bit
of the current block. A seven-state binary decision tree reads
its three bits. For each possible symbol, fewer than ten further
states return to the first bit, write the new three-bit block and
move to the first bit of the adjacent block. Thus fewer than 128
binary states per simulated state suffice. This simulation works
on the initially blank tape and requires no end-of-tape convention.
More explicitly, number the current block's cells 0,1,2. The
three read transitions move right, right, left, leaving the head
at cell 1 with the read word in finite control. Preserve cell 1
and move left, then write the three new bits while moving right,
right, and the simulated direction. For a right move the head is
already at cell 3. For a left move it is at cell 1; four
symbol-preserving left moves reach cell -3. Thus at most eight
post-read states per read word suffice, with no stationary move.
A halting transition enters the halting state at the end of the
same write-and-move sequence. Every other transition returns
block-aligned to the next state's read routine.
In particular the generous bound
\[
                    2^{20}(L+1)(R+1)
\]
counts all working states of the resulting binary machine.

We hardwire finite input without assuming a decoder. Starting at
zero, an input counter is replaced by twice itself plus the next
bit, reading the fixed bit string from most significant to least
significant bit. Copy followed by move doubles the counter, and
one additional increment writes a one bit. Fewer than sixteen
counter instructions per bit, and two reusable scratch counters,
suffice. This is a finite chain of instructions, not a variable
input tape. The literal width 1000 in the coefficient interpreter
costs at most 1001 additional instructions.
For this input prefix specifically, every copy uses the same two
scratch counters: copy the input counter to the first while using
the second to restore it, then move the first into the input
counter. Both scratch counters are zero afterwards. The next bit
may therefore reuse them. Fresh labels are generated for each bit
but fresh counters are not. This explicit reuse is distinct from
the convenient fresh-per-node allocation for the fixed core.

For the record search, $0\le j<U=2^{2^{600}}$ requires at most
$2^{600}$ bits, including the representation 0 when $j=0$.
After hardwiring this input the counter program has
$L<2^{605}$ and $R<2^{187}$. Its binary machine consequently has
fewer than $2^{20+606+188}=2^{814}<2^{900}$ working states, in particular fewer than
$2^{1000}$. The final loop has $G(r)$ increments, each taking at
least one binary transition, and initialization and termination
take further transitions. Its total halting time is strictly
greater than $G(r)$.

For the coefficient machine, put $s=2^{1000}$. A complete table
has $\ell=2006s$ bits, including the written bit, direction bit,
halting bit and 1000-bit next-state index for each of $2s$
entries. Hardwiring all these bits and width 1000 gives
$L<2^{1016}$ and $R<2^{149}$. The same construction therefore
uses fewer than $2^{1200}$ binary working states, in particular
fewer than $2^{2000}$. No input decoder or small universal
machine is assumed in either count. Unused states pad these
machines to exactly the state counts occurring in $\BB$.

These are constructive upper bounds, not claims that an enormous
transition table has been printed. The finite syntax checks verify
the stated source restrictions. The compilation templates and
tape simulation above provide the state bound; interpreter test
cases alone would not provide it.

For the actual finite input $j=j_*$, the record program halts and its
running time is strictly larger than $G(r_*)$. Its compiled state count
and padding therefore give
\[
 G(r_*)<\BB(2^{1000}).
\]
The existence of this finite hardwired input does not require an
algorithm deciding whether a record is the final record.


Taking $N_*=G(r_*)$ proves Theorem~\ref{q198bb:threshold}.
Its all-larger-parameter form applies at $N_{\rm BB}=\BB(2^{1000})$.
The independent coefficient crossover check is given next.


\subsection{A strict Busy Beaver bound for the coefficients}
\begin{lemma}\label{q198bb:coefficient}
For the fixed binary, two-way-tape, running-time convention above,
\[
 149^{4\BB(2^{1000})+2}+1<\BB(2^{2000}).
\]
\end{lemma}
\begin{proof}
Put $s=2^{1000}$. There are finitely many $s$-state transition
tables and at least one halting table. Choose a table whose
blank-tape halting time is $T=\BB(s)$. This is an existence
choice, not an algorithm for finding a maximizing machine.
Number the working states $0,\ldots,s-1$, with state 0 initial.
For each pair of working state and scanned symbol, record the
written bit, direction bit, halting bit and 1000 bits for the next
working-state index. The index is ignored on a halting transition.
The complete table therefore uses exactly
\[
                        \ell=2006\cdot2^{1000}
\]
bits, allowing leading zero bits.

The displayed fixed integer interpreter simulates this table.
Its left and right integers store tape symbols with the nearest
cell as the least significant bit, and its separate symbol
register stores the cell under the head. Writing and then moving
pushes the written bit onto one integer and pops the other.
Trailing zero bits represent arbitrarily many blank cells, so
this describes the full two-way infinite tape. The loop counts
the final transition, including its write and move, before
observing the halt bit. It therefore returns exactly $T$.

The program next computes $H=149^{4T+2}+1$ and performs $H+1$
increments before halting. Section~\ref{q198bb:compiler} constructs
a binary machine with fewer than $2^{1200}<2^{2000}$ working
states for this fixed interpreter and its hardwired table.
Its final loop alone takes more than $H$ transitions. Padding
with unused working states gives a machine of exactly
$2^{2000}$ states without changing the run. Thus
$H<\BB(2^{2000})$.

Finally a chain of 54 working states, each moving right and
passing to the next state, with the last transition halting,
runs for 54 transitions on the blank tape. Padding it proves
$\BB(2^{1000})\ge54>53$. Hence the separately checked coefficient
formula $H(P_N)=\max\{H_0,149^{4N+2}\}=149^{4N+2}$ for $N\ge53$
applies independently of the arithmetic-threshold argument. Therefore
\[
 H(P_{\rm BB})=149^{4N_{\rm BB}+2}
 <149^{4N_{\rm BB}+2}+1<\BB(2^{2000}).
\]
\end{proof}


Theorem~\ref{q198bb:threshold}, Lemma~\ref{q198bb:coefficient}, and the exact
coefficient formula in Section~\ref{q198bb:inputs} prove
Theorem~\ref{q198bb:fixed}. The construction now selects one polynomial
using the specified Busy Beaver function. No assertion here computes
that function or gives a decimal evaluation of the selected parameter.


\subsection{Finite algebra for the fixed embedding}\label{q198bb:finite-algebra}
For $Q=(x,y)$ and $R=(u,v)$ on $E_0$, put $\delta=u-x$.
Let $s=(x(Q+R),y(Q+R),1)$ and
$t=(x(Q-R),y(Q-R),1)$ wherever these affine expressions are defined.
Then $B_{ij}=\delta^3s_it_j$ reduces, using $y^2=x^3-x$ and
$v^2=u^3-u$, to the following polynomials. Each has degree at most
two in $(x,y)$ and at most two in $(u,v)$, so it homogenizes to the
claimed bidegree $(2,2)$ section.
\[
\begin{aligned}
B_{1,1}={}&u-u^{2}y^{2}-x\\
&+xu^{2}+x^{2}v^{2}-x^{2}u
\end{aligned}
\]

\[
\begin{aligned}
B_{1,2}={}&v+y+uy^{2}v\\
&+u^{2}y+xyv^{2}+2xuv\\
&+2xuy+x^{2}v
\end{aligned}
\]

\[
\begin{aligned}
B_{1,3}={}&-2uyv-uy^{2}-u^{2}\\
&+xv^{2}+2xyv+x^{2}
\end{aligned}
\]

\[
\begin{aligned}
B_{2,1}={}&-v+y-uy^{2}v\\
&+u^{2}y+xyv^{2}-2xuv\\
&+2xuy-x^{2}v
\end{aligned}
\]

\[
\begin{aligned}
B_{2,2}={}&1+y^{2}v^{2}-u^{2}\\
&-4xu-x^{2}-3x^{2}u^{2}
\end{aligned}
\]

\[
\begin{aligned}
B_{2,3}={}&yv^{2}-y^{2}v+uv\\
&-2uy+2xv-xy\\
&+3xu^{2}y-3x^{2}uv
\end{aligned}
\]

\[
\begin{aligned}
B_{3,1}={}&2uyv-uy^{2}-u^{2}\\
&+xv^{2}-2xyv+x^{2}
\end{aligned}
\]

\[
\begin{aligned}
B_{3,2}={}&yv^{2}+y^{2}v-uv\\
&-2uy-2xv-xy\\
&+3xu^{2}y+3x^{2}uv
\end{aligned}
\]

\[
\begin{aligned}
B_{3,3}={}&v^{2}-y^{2}+u\\
&-x-3xu^{2}+3x^{2}u
\end{aligned}
\]

The formulas can be checked by the ordinary chord-and-tangent formulas,
clearing powers of $u-x$, and reducing by the two curve equations.
Their identities on this dense affine open determine their global
sections. The theorem of the square gives the line-bundle identification
used for the addition-subtraction morphism. Its trivialization here is
multiplication by $(u-x)^3$.



\subsection{The record-search program}\label{q198bb:programs}
The following is executable notation for unbounded integer semantics.
The finite input is $j$. Only the function definitions are displayed;
reading them does not evaluate the production constant. Line wrapping
in three expressions does not change the program's parsed syntax.
\begingroup\small
\begin{verbatim}
"""Reference program with mathematical, unbounded integer semantics.

The production constant is deliberately not evaluated when this file is imported.
The proof uses only fixed-size tuples, acyclic function calls, and integer loops.
Python is convenient executable notation for that register-machine program.
"""


def gcd(a, b):
    if a < 0:
        a = -a
    if b < 0:
        b = -b
    while b != 0:
        remainder = a % b
        a = b
        b = remainder
    return a


def field(a, b, denominator):
    if denominator < 0:
        a = -a
        b = -b
        denominator = -denominator
    divisor = gcd(gcd(a, b), denominator)
    return (a // divisor, b // divisor, denominator // divisor)


def field_add(left, right):
    a, b, c = left
    d, e, f = right
    return field(a*f + d*c, b*f + e*c, c*f)


def field_neg(value):
    a, b, c = value
    return (-a, -b, c)


def field_sub(left, right):
    return field_add(left, field_neg(right))


def field_mul(left, right):
    a, b, c = left
    d, e, f = right
    return field(a*d + 2*b*e, a*e + b*d, c*f)


def field_div(left, right):
    a, b, c = right
    reciprocal = field(c*a, -c*b, a*a - 2*b*b)
    return field_mul(left, reciprocal)


def point_add(left, right):
    if left is None:
        return right
    if right is None:
        return left
    x1, y1 = left
    x2, y2 = right
    if x1 == x2:
        if y1 != y2 or y1 == (0, 0, 1):
            return None
        numerator = field_sub(
            field_mul((3, 0, 1), field_mul(x1, x1)), (11, -6, 1))
        slope = field_div(numerator, field_mul((2, 0, 1), y1))
    else:
        slope = field_div(field_sub(y2, y1), field_sub(x2, x1))
    x3 = field_sub(field_sub(field_mul(slope, slope), x1), x2)
    y3 = field_sub(field_mul(slope, field_sub(x1, x3)), y1)
    return (x3, y3)


def multiple(n):
    point = ((1, 1, 1), (0, 2, 1))
    result = None
    remaining = n
    while remaining > 0:
        if remaining % 2 != 0:
            result = point_add(result, point)
        remaining = remaining // 2
        if remaining > 0:
            point = point_add(point, point)
    return result


def point_encoded(q):
    # Positive q encodes m=q//4 and one of four rational two-torsion points.
    # q=0 alone is omitted, so every encoded point is finite.
    m = q // 4
    torsion = q % 4
    point = multiple(m)
    if torsion == 0:
        translate = None
    elif torsion == 1:
        translate = ((0, 0, 1), (0, 0, 1))
    elif torsion == 2:
        translate = ((3, -1, 1), (0, 0, 1))
    else:
        translate = ((-3, 1, 1), (0, 0, 1))
    return point_add(point, translate)

def rational_add(left, right):
    a, b = left
    c, d = right
    numerator = a*d + c*b
    denominator = b*d
    divisor = gcd(numerator, denominator)
    return (numerator // divisor, denominator // divisor)


def trace_height_bits(n1, n2, n3, n4):
    # T = 7 plus one quarter of the trace, hence 7 plus half the sum
    # of the rational coefficients of the four x-coordinates.
    x1, y1 = point_encoded(n1)
    x2, y2 = point_encoded(n2)
    x3, y3 = point_encoded(n3)
    x4, y4 = point_encoded(n4)
    numerator, denominator = rational_add(
        rational_add((x1[0], x1[2]), (x2[0], x2[2])),
        rational_add((x3[0], x3[2]), (x4[0], x4[2])))
    numerator = numerator + 14*denominator
    denominator = 2*denominator
    divisor = gcd(numerator, denominator)
    numerator = numerator // divisor
    denominator = denominator // divisor
    if numerator < 0:
        numerator = -numerator
    height = denominator
    if numerator > height:
        height = numerator
    bits = 0
    while height > 0:
        height = height // 2
        bits = bits + 1
    return bits


def power(base, exponent):
    result = 1
    remaining = exponent
    while remaining > 0:
        result = result*base
        remaining = remaining - 1
    return result


def threshold_after_record(j, scale):
    record = scale
    count = 0
    radius = 0
    while count < j:
        radius = radius + 1
        n1 = 1
        while n1 <= radius and count < j:
            n2 = 1
            while n2 <= radius and count < j:
                n3 = 1
                while n3 <= radius and count < j:
                    n4 = 1
                    while n4 <= radius and count < j:
                        if (n1 == radius or n2 == radius
                                or n3 == radius or n4 == radius):
                            size = ((n1//4)*(n1//4)
                                    + (n2//4)*(n2//4)
                                    + (n3//4)*(n3//4)
                                    + (n4//4)*(n4//4))
                            if size > scale*record:
                                bits = trace_height_bits(n1, n2, n3, n4)
                                if size > scale*bits:
                                    record = size
                                    count = count + 1
                        n4 = n4 + 1
                    n3 = n3 + 1
                n2 = n2 + 1
            n1 = n1 + 1
    return power(scale, 32)*(record + 1)


def machine(j):
    scale = power(2, power(2, 600))
    threshold = threshold_after_record(j, scale)
    clock = 0
    while clock < threshold:
        clock = clock + 1
    return threshold
\end{verbatim}
\endgroup

\subsection{The coefficient-bound program}\label{q198bb:coefficient-program}
The following fixed program takes the encoded transition table and
\texttt{width}$=1000$. Its loop counts every simulated transition,
including the transition that sets the halting bit. The final loop
performs one more increment than the computed coefficient height.
\begingroup\small
\begin{verbatim}
"""A fixed integer-register interpreter for the coefficient BB bound.

For width w, states are 0,...,2**w-1. Each transition occupies w+3 bits:
write bit, direction bit (0=left), halt bit, and next-state bits.
The record for (state, symbol) begins at bit (w+3)*(2*state+symbol).
All arithmetic has unbounded integer semantics. No BB value is computed.
"""


def pow2(exponent):
    result = 1
    while exponent > 0:
        result = result * 2
        exponent = exponent - 1
    return result


def simulate(table, width):
    left = 0
    right = 0
    symbol = 0
    state = 0
    ticks = 0
    halted = 0
    while halted == 0:
        offset = (width + 3) * (2 * state + symbol)
        record = (table // pow2(offset)) % pow2(width + 3)
        written = record % 2
        direction = (record // 2) % 2
        halted = (record // 4) % 2
        state = record // 8
        if direction == 0:
            right = 2 * right + written
            symbol = left % 2
            left = left // 2
        else:
            left = 2 * left + written
            symbol = right % 2
            right = right // 2
        ticks = ticks + 1
    return ticks


def machine(table, width):
    ticks = simulate(table, width)
    exponent = 4 * ticks + 2
    value = 1
    while exponent > 0:
        value = value * 149
        exponent = exponent - 1
    bound = value + 1
    clock = 0
    while clock <= bound:
        clock = clock + 1
    return bound
\end{verbatim}
\endgroup
